%---
% title: "Tôi đã giải Toán THPT bằng “tà đạo” — bài 2: Một phép chiếu thay sáu công thức Oxyz"
% date: 2026-08-06
% tags: [dai-so-tuyen-tinh, thpt, oxyz, hinh-hoc]
% summary: "Bài 2 chuỗi «Tôi đã giải Toán THPT bằng phương pháp tà đạo»: cách tôi dùng đúng một công thức hình chiếu proj_u(v) = ((u·v)/(u·u))u của đại số tuyến tính để xử gọn hình chiếu, chân đường vuông góc và các dạng khoảng cách trong Oxyz — kèm chứng minh, ví dụ từ đề chính thức và bài tập có đáp án."
%---
\documentclass[12pt,a4paper]{article}

% Khối khai báo dưới đây CHỈ dùng khi biên dịch PDF trực tiếp
% (Overleaf, pdfLaTeX). Blog không cần nó: Pandoc tự bỏ qua phần này.
\usepackage[utf8]{inputenc}
\usepackage[T5]{fontenc}
\usepackage[vietnamese]{babel}
\usepackage{amsmath,amssymb}
\usepackage[a4paper,margin=2.2cm]{geometry}
\usepackage{tikz}
\usetikzlibrary{arrows.meta}

% Hình minh họa chỉ dành cho bản web: pdfLaTeX bỏ trống,
% Pandoc (qua build script) gỡ wrapper và xuất thẻ <img>.
\newcommand{\blogfig}[1]{}

\begin{document}

\section{Nếu chỉ được mang một thứ vào phòng thi Oxyz}

Nếu được chọn đúng một vũ khí duy nhất cho phần hình học không gian, tôi chọn \textbf{phép chiếu} --- một khái niệm của đại số tuyến tính mà tôi gọi nôm na là phép ``bóng đổ''.

Sáu dạng bài tôi từng học thuộc riêng lẻ gồm: hình chiếu lên đường thẳng, chân đường vuông góc, khoảng cách điểm--đường thẳng, hình chiếu vectơ lên mặt phẳng, khoảng cách điểm--mặt phẳng và vị trí tương đối của mặt cầu. Tất cả đều quy về cùng một phép tách. Bài trước đã giới thiệu tích vô hướng. Bài này dùng nó để dọn sạch sáu dạng trên.

Ký hiệu giữ nguyên quy ước bài trước: vectơ tự do in đậm ($\mathbf{u}, \mathbf{v}, \mathbf{w}, \mathbf{n}$), vectơ nối điểm có mũi tên ($\overrightarrow{AM} = \overrightarrow{OM}-\overrightarrow{OA}$).

\section{Hành trang tôi cần}

Từ bài tích vô hướng, tôi mang theo bốn món:

\begin{itemize}
  \item $\mathbf{u}\cdot\mathbf{v} = u_1v_1+u_2v_2+u_3v_3 = \|\mathbf{u}\|\|\mathbf{v}\|\cos\theta$.
  \item Trực giao: $\mathbf{u} \perp \mathbf{v} \Leftrightarrow \mathbf{u}\cdot\mathbf{v} = 0$.
  \item Phân phối: $(\mathbf{a}+\mathbf{b})\cdot\mathbf{c} = \mathbf{a}\cdot\mathbf{c} + \mathbf{b}\cdot\mathbf{c}$ và $(t\,\mathbf{a})\cdot\mathbf{b} = t(\mathbf{a}\cdot\mathbf{b})$.
  \item Chuẩn: $\|\mathbf{v}\| = \sqrt{\mathbf{v}\cdot\mathbf{v}}$; khoảng cách $AB = \|\overrightarrow{OB}-\overrightarrow{OA}\|$.
\end{itemize}

Và hai khái niệm Oxyz quen thuộc: đường thẳng $d$ qua $A$ với vector chỉ phương $\mathbf{u}$ (mọi điểm trên $d$ có dạng $A+t\,\mathbf{u}$); mặt phẳng $(P): ax+by+cz+d=0$ với pháp vector $\mathbf{n}=(a;b;c)$ đứng vuông góc với mọi vectơ nằm trong mặt.

\section{Trực giác: cái bóng biết nói}

Một chiếc gậy trên mặt đất, mặt trời chiếu thẳng đứng, một con chim lơ lửng giữa trời. Bóng của chim rơi vào đâu? Ba điều hiển nhiên:

\begin{enumerate}
  \item Bóng nằm trên mặt đất.
  \item Đoạn nối chim với bóng thẳng đứng --- vuông góc với đất.
  \item Bóng là điểm gần chim nhất trên mặt đất.
\end{enumerate}

Cái hay là điều thứ ba không cần phải tin mù, vì nó là hệ quả của điều thứ hai. Gọi $H$ là chân vuông góc, lấy bất kỳ điểm $X$ khác trên mặt đất, Pythagoras cho:

$$
MX^2 = MH^2 + HX^2 > MH^2
$$

Dấu bằng chỉ xảy ra khi $X=H$. Nghĩa là ``vuông góc'' và ``gần nhất'' là một tính chất nhìn từ hai phía. Khi tôi nhận ra điều này, mọi công thức khoảng cách bắt đầu mất đi tính thần kỳ.

\begin{center}
\begin{tikzpicture}[>=Stealth,line cap=round,scale=1]
  \draw[line width=1.2pt,gray!80] (-0.5,0) -- (8.5,0);
  \node[gray!80,below,font=\small] at (8.2,0) {$\Delta$};
  \shade[ball color=amber!85!orange] (7.6,3.4) circle (0.32);
  \foreach \a in {0,45,...,315}
    \draw[amber!80!orange,thick] ([shift={(\a:0.45)}]7.6,3.4) -- ([shift={(\a:0.64)}]7.6,3.4);
  \draw[dashed,amber!70!black,thick,->] (7.3,3.22) -- (5.75,2.05);
  \fill (5.6,1.9) circle (2.2pt);
  \node[above left,font=\small] at (5.6,1.9) {chim $M$};
  \draw[dashed,purple!80!black,thick] (5.6,1.9) -- (5.6,0);
  \fill[purple!80!black] (5.6,0) circle (2.4pt);
  \node[purple!80!black,below,font=\small] at (5.6,-0.05) {hình chiếu $H$};
\end{tikzpicture}
\end{center}

\blogfig{\includegraphics{/images/la-projection/bong-do.svg}}

\textbf{Hình 1.} Bóng của $M$ trên $\Delta$: nằm trên $\Delta$, vuông góc, và gần $M$ nhất --- ba tính chất là một.

\section{Công thức tôi xây, từng bước có kiểm chứng}

Cho $\Delta$ qua gốc $O$, vector chỉ phương $\mathbf{u} \neq \mathbf{0}$, và một vectơ $\mathbf{v}$. Tôi muốn hình chiếu của $\mathbf{v}$ lên $\Delta$.

\textbf{Bước 1.} Mọi điểm trên $\Delta$ có dạng $t\,\mathbf{u}$ ($t$ thực). Vậy bài toán thu thành: tìm số $t$ đẹp nhất.

\textbf{Bước 2.} Điều kiện ``đẹp nhất'' theo trực giác bóng đổ: $(\mathbf{v} - t\,\mathbf{u})\cdot\mathbf{u} = 0$ --- nếu đoạn nối còn xiên thì vẫn trượt dọc $\Delta$ được thêm chút nữa để gần hơn.

\textbf{Bước 3.} Khai triển bằng tính phân phối:

$$
\mathbf{v}\cdot\mathbf{u} - t(\mathbf{u}\cdot\mathbf{u}) = 0
\quad\Longrightarrow\quad
t = \frac{\mathbf{v}\cdot\mathbf{u}}{\mathbf{u}\cdot\mathbf{u}}
$$

Chia an toàn vì $\mathbf{u}\neq\mathbf{0}$ kéo theo $\mathbf{u}\cdot\mathbf{u}=\|\mathbf{u}\|^2>0$. Vậy:

$$
\boxed{\;\operatorname{proj}_{\mathbf{u}}(\mathbf{v}) = \frac{\mathbf{v} \cdot \mathbf{u}}{\mathbf{u} \cdot \mathbf{u}}\, \mathbf{u}
= \frac{u_1 v_1 + u_2 v_2 + u_3 v_3}{u_1^2 + u_2^2 + u_3^2}\;(u_1;u_2;u_3)\;}
$$

\subsection*{Tôi kiểm tra lại công thức làm đúng đúng hai việc}

Thứ nhất, nó nằm dọc $\mathbf{u}$ (có dạng $t\,\mathbf{u}$). Thứ hai, phần dư vuông góc:

$$
\bigl(\mathbf{v} - \operatorname{proj}_{\mathbf{u}}(\mathbf{v})\bigr)\cdot\mathbf{u}
= \mathbf{v}\cdot\mathbf{u} - \frac{\mathbf{v}\cdot\mathbf{u}}{\mathbf{u}\cdot\mathbf{u}}\,(\mathbf{u}\cdot\mathbf{u})
= 0 \;\checkmark
$$

Hai việc cần của hình chiếu đều đạt.

\subsection*{Và đây là lúc chữ ``tuyến tính'' xuất hiện}

Điều khiến tôi chắc chắn khái niệm này thuộc về đại số tuyến tính: phép gán $\mathbf{v} \mapsto \operatorname{proj}_{\mathbf{u}}(\mathbf{v})$ là một \textbf{ánh xạ tuyến tính}, tức bảo toàn tổ hợp tuyến tính:

$$
\operatorname{proj}_{\mathbf{u}}(a\,\mathbf{v} + b\,\mathbf{w})
= a\,\operatorname{proj}_{\mathbf{u}}(\mathbf{v}) + b\,\operatorname{proj}_{\mathbf{u}}(\mathbf{w})
$$

Kiểm tra chỉ mất một dòng: tử số $(a\mathbf{v}+b\mathbf{w})\cdot\mathbf{u}$ phân phối thành $a(\mathbf{v}\cdot\mathbf{u})+b(\mathbf{w}\cdot\mathbf{u})$, hằng số kéo ra ngoài. Chính cấu trúc này làm phép chiếu trở nên quan trọng ở đại học (trực giao hóa Gram--Schmidt, hồi quy tuyến tính), chứ không chỉ trong phòng thi.

Ba phép thử quen thuộc cũng đạt: chiếu lên trục Ox ra $(v_1;0;0)$; chiếu vectơ đã nằm dọc $\mathbf{u}$ ra chính nó; vectơ vuông góc thì bóng sụp về $\mathbf{0}$.

\section{Tách vectơ: phần nằm trên, phần thẳng đứng}

$$
\mathbf{v} = \underbrace{\operatorname{proj}_{\mathbf{u}}(\mathbf{v})}_{\text{nằm dọc } \mathbf{u}}
\;+\;
\underbrace{\mathbf{v}_{\perp}}_{\perp\, \mathbf{u}},
\qquad
\mathbf{v}_{\perp} = \mathbf{v} - \operatorname{proj}_{\mathbf{u}}(\mathbf{v})
$$

Một phát ăn ngay hai dạng câu hỏi: muốn hình chiếu lấy vế đầu, muốn khoảng cách lấy độ dài vế sau $\|\mathbf{v}_{\perp}\|$ (đúng ngắn nhất theo lập luận Pythagoras ở mục bóng đổ).

\begin{center}
\begin{tikzpicture}[>=Stealth,line cap=round,scale=1.05]
  \draw[thick,gray!70] (-0.8,-0.4) -- (6.4,2.8);
  \node[gray!70,right,font=\small] at (6.4,2.8) {$\Delta$};
  \coordinate (O) at (0,0);
  \coordinate (H) at (4,2);
  \coordinate (M) at (4,3.5);
  \draw[->,line width=1.3pt,blue!70!black] (O) -- (H);
  \node[blue!70!black,above,sloped,font=\small] at (2,1) {$\operatorname{proj}_{\mathbf{u}}(\mathbf{v})$};
  \draw[->,line width=1.1pt,red!70!black] (O) -- (M);
  \node[red!70!black,left,font=\small] at (4,3.5) {$\mathbf{v}$};
  \draw[dashed,line width=1.1pt,purple!80!black] (H) -- (M);
  \node[purple!80!black,right,font=\small] at (4,2.75) {$\mathbf{v}_{\perp}$};
  \draw[->,line width=1.3pt,teal!70!black] (0.5,0.25) -- (1.9,0.95);
  \node[teal!70!black,above,font=\small] at (1.9,0.95) {$\mathbf{u}$};
  \fill (O) circle (1.6pt); \node[below left,font=\small] at (O) {$O$};
  \fill[blue!70!black] (H) circle (1.6pt); \node[below right,font=\small] at (H) {$H$};
  \fill[red!70!black] (M) circle (1.6pt); \node[right,font=\small] at (4.15,3.55) {$M$};
  \draw[purple!80!black] (3.73,1.87) -- (3.73,2.17) -- (4,2.32);
\end{tikzpicture}
\end{center}

\blogfig{\includegraphics{/images/la-projection/phan-tach-hinh-chieu.svg}}

\textbf{Hình 2.} $\mathbf{v} = \operatorname{proj}_{\mathbf{u}}(\mathbf{v}) + \mathbf{v}_{\perp}$: phần nằm trên $\Delta$ và phần thẳng đứng với $\Delta$.

\section{Sáu dạng bài, một phép tách}

\emph{Ví dụ có ghi rõ nguồn là bài chính thức; các ví dụ còn lại tôi biên soạn theo cấu trúc và mức độ của đề tốt nghiệp THPT.}

\subsection*{Dạng 1. Khoảng cách từ điểm đến đường thẳng}

Quy trình tôi đi mãi thành phản xạ: lấy $\mathbf{v} = \overrightarrow{AM} = \overrightarrow{OM} - \overrightarrow{OA}$ (mẹo mượn điểm gốc để quay về tình huống đã giải), chiếu lên $\mathbf{u}$, rồi:

$$
d(M,d) = \bigl\| \mathbf{v} - \operatorname{proj}_{\mathbf{u}}(\mathbf{v}) \bigr\|
$$

\textbf{Ví dụ 1.} Cho $d:\;\dfrac{x-1}{2}=y=\dfrac{z-1}{2}$ và $M(5;3;-3)$. Tính khoảng cách từ $M$ đến $d$.

\textbf{Lời giải.} Đường thẳng qua $A(1;0;1)$, chỉ phương $\mathbf{u}=(2;1;2)$, vector nối $\mathbf{v} = \overrightarrow{AM} = (4;3;-4)$. Tôi tính được $\mathbf{v}\cdot\mathbf{u}=3$, $\mathbf{u}\cdot\mathbf{u}=9$, nên $t=\frac{1}{3}$:

$$
\operatorname{proj}_{\mathbf{u}}(\mathbf{v}) = \left(\frac{2}{3};\,\frac{1}{3};\,\frac{2}{3}\right),
\qquad
\mathbf{v}_{\perp} = \left(\frac{10}{3};\,\frac{8}{3};\,-\frac{14}{3}\right)
$$

$$
d(M,d) = \|\mathbf{v}_{\perp}\| = \sqrt{\frac{100+64+196}{9}} = 2\sqrt{10}
$$

Chân đường vuông góc rơi ra luôn: $H = A + \operatorname{proj}_{\mathbf{u}}(\mathbf{v}) = \left(\dfrac{5}{3};\,\dfrac{1}{3};\,\dfrac{5}{3}\right)$.

\subsection*{Dạng 2. Hình chiếu của điểm lên đường thẳng}

$$
\boxed{\;H = A + \operatorname{proj}_{\mathbf{u}}(\overrightarrow{AM})\;}
$$

Tại sao phải cộng $A$? Vì hình chiếu tính ra là vectơ kể từ gốc tọa độ; muốn đặt nó lên đúng đường thẳng thì phải ``ghim'' vào điểm gốc $A$. Chứng minh đủ điều kiện tôi đã trình bày ở phần công thức: $H-A$ nằm dọc $\mathbf{u}$ (nên $H \in d$) và $(M-H)\cdot\mathbf{u}=0$ (nên $MH \perp d$).

\textbf{Ví dụ 2.} Cho $d:\;\dfrac{x-2}{1}=\dfrac{y-1}{2}=\dfrac{z+1}{-2}$ và $M(5;1;-1)$.

Đường thẳng qua $A(2;1;-1)$, chỉ phương $\mathbf{u}=(1;2;-2)$, $\mathbf{v}=\overrightarrow{AM}=(3;0;0)$:

$$
t = \frac{3}{9} = \frac{1}{3}
\quad\Longrightarrow\quad
H = A + \operatorname{proj}_{\mathbf{u}}(\mathbf{v}) = \left(\frac{7}{3};\,\frac{5}{3};\,-\frac{5}{3}\right)
$$

Kiểm tra: $M-H = \left(\dfrac{8}{3};-\dfrac{2}{3};\dfrac{2}{3}\right)$, nhân với $\mathbf{u}$ ra $0$ --- vuông góc chuẩn. Khoảng cách: $d(M,d)=\sqrt{\frac{64+4+4}{9}}=2\sqrt{2}$.

Một phép chiếu trả về cả hai đáp án mà đề hay tách làm hai câu.

\subsection*{Dạng 3. Hình chiếu của một vectơ lên mặt phẳng}

Đến mặt phẳng, ban đầu tôi tưởng sẽ khó hơn đường thẳng. Thực ra nhờ pháp vector, nó lại còn gọn hơn:

$$
\boxed{\;\mathbf{w}_{\parallel} = \mathbf{w} - \operatorname{proj}_{\mathbf{n}}(\mathbf{w})\;}
$$

\textbf{Vì sao đúng?} Chỉ cần kiểm tra $\mathbf{w}_{\parallel}$ trực giao với $\mathbf{n}$:

$$
\mathbf{w}_{\parallel}\cdot\mathbf{n}
= \mathbf{w}\cdot\mathbf{n} - \frac{\mathbf{w}\cdot\mathbf{n}}{\mathbf{n}\cdot\mathbf{n}}\,(\mathbf{n}\cdot\mathbf{n}) = 0
$$

Trực giao với $\mathbf{n}$ nghĩa là nằm trong mặt phẳng qua gốc có pháp vector $\mathbf{n}$. Pháp vector lấy từ hệ số phương trình ($ax+by+cz+d=0$ cho ngay $\mathbf{n}=(a;b;c)$), hoặc từ tích có hướng nếu đề chỉ cho hai vector chỉ phương.

\textbf{Ví dụ 3 (làm quen).} Chiếu $\mathbf{w}=(1;2;3)$ lên mặt Oxy ($\mathbf{n}=(0;0;1)$): $\operatorname{proj}_{\mathbf{n}}(\mathbf{w})=(0;0;3)$, nên $\mathbf{w}_{\parallel}=(1;2;0)$. Bóng trên sàn Oxy, đúng như mắt thường thấy.

\subsection*{Dạng 4. Khoảng cách từ điểm đến mặt phẳng}

$$
d\bigl(M,(P)\bigr) = \frac{|\overrightarrow{AM}\cdot\mathbf{n}|}{\|\mathbf{n}\|}
$$

Công thức nổi tiếng của sách giáo khoa suy ra trong vài dòng: chọn $A$ bất kỳ thuộc $(P)$ thì $ax_A+by_A+cz_A+d=0$, tức $\overrightarrow{OA}\cdot\mathbf{n}=-d$, nên:

$$
\overrightarrow{AM}\cdot\mathbf{n}
= (\overrightarrow{OM}-\overrightarrow{OA})\cdot\mathbf{n}
= ax_0+by_0+cz_0+d
$$

$$
\Longrightarrow\quad
d\bigl(M,(P)\bigr) = \frac{|ax_0 + by_0 + cz_0 + d|}{\sqrt{a^2+b^2+c^2}}
$$

Lần đầu tự suy ra được công thức này, tôi ngừng học thuộc nó luôn.

\textbf{Ví dụ 4.} Cho $(P): 2x-2y+z-10=0$ và $M(1;2;3)$.

\emph{Cách công thức}: $|2-4+3-10|/\sqrt{9} = 9/3 = 3$. \emph{Cách hình chiếu}: chọn $A(5;0;0)\in(P)$, $\overrightarrow{AM}=(-4;2;3)$, $\overrightarrow{AM}\cdot\mathbf{n} = -9$, ra cùng $3$. Hai đường một đáp án.

Hệ quả tôi hay dùng: đường thẳng $\Delta$ song song $(P)$ thì mọi điểm của nó cách đều, nên lấy điểm bất kỳ trên $\Delta$ rồi tính Dạng 4.

\subsection*{Dạng 5. Hình chiếu của một điểm lên mặt phẳng}

$$
\boxed{\;H = M - \operatorname{proj}_{\mathbf{n}}(\overrightarrow{AM})\;}
$$

Công thức này tôi từng chép máy móc. Sau khi kiểm tra, nó thuyết phục hẳn: $M-H$ nằm dọc $\mathbf{n}$ (nên $MH \perp (P)$), và $(H-A)\cdot\mathbf{n} = \overrightarrow{AM}\cdot\mathbf{n} - \overrightarrow{AM}\cdot\mathbf{n} = 0$ (nên $H$ nằm đúng trong $(P)$). Cả hai điều kiện của hình chiếu đều đạt.

\textbf{Ví dụ 5.} Cho $(P): 2x+2y+z-9=0$ và $M(3;2;1)$. Chọn $A(4;0;1)\in(P)$, $\overrightarrow{AM}=(-1;2;0)$:

$$
\overrightarrow{AM}\cdot\mathbf{n} = -2+4+0 = 2,
\qquad
H = M - \frac{2}{9}\,(2;2;1) = \left(\frac{23}{9};\,\frac{14}{9};\,\frac{7}{9}\right)
$$

Tôi luôn kiểm tra hai điều kiện với đáp án phân số: thay $H$ vào $(P)$ được $0$; và $M-H = \frac{2}{9}\mathbf{n}$ cùng phương pháp vector. Khoảng cách rơi ra: $\|M-H\| = \frac{2}{3}$.

\subsection*{Dạng 6 (mở rộng). Mặt cầu}

Khoảng cách từ tâm $I$ đến mặt phẳng --- tiêu chí biện luận vị trí tương đối --- chỉ là Dạng 4 đổi nhân vật. Nếu cắt nhau, bán kính đường tròn giao tuyến là cạnh còn lại của tam giác vuông: $r = \sqrt{R^2-d^2}$.

\textbf{Ví dụ 6 (lấy nguyên văn từ: Thi THPT Quốc gia 2019 --- Mã đề 101 --- Câu 21).} Cho mặt cầu $(S): x^2+y^2+z^2-2x-2z+7=0$. Bán kính bằng

\textbf{A.} $\sqrt{7}$. \qquad \textbf{B.} $9$. \qquad \textbf{C.} $3$. \qquad \textbf{D.} $\sqrt{15}$.

Hoàn thành bình phương: $(x-1)^2+y^2+(z-1)^2 = 9$, vậy $R=3$, \textbf{chọn C}. Một câu 30 giây nếu nhớ công thức, nhưng giờ tôi biết con số 3 đó đến từ đâu.

\textbf{Ví dụ 7 (tổng hợp).} Cho $(S): x^2+y^2+z^2-2x-4y+2z-3=0$ và $(P): x+2y+2z-6=0$. Tâm $I(1;2;-1)$, $R=\sqrt{9}=3$; khoảng cách $d(I,(P)) = |1+4-2-6|/3 = 1 < R$: mặt phẳng \textbf{cắt} mặt cầu, giao tuyến có bán kính $r=\sqrt{8}=2\sqrt{2}$.

\section{Bảng công thức tôi mang theo}

\begin{center}
\begin{tabular}{|l|l|}
\hline
\textbf{Yêu cầu} & \textbf{Công thức} \\
\hline
Hình chiếu lên phương $\mathbf{u}$ & $\operatorname{proj}_{\mathbf{u}}(\mathbf{v}) = \dfrac{\mathbf{v}\cdot\mathbf{u}}{\mathbf{u}\cdot\mathbf{u}}\,\mathbf{u}$ \\
\hline
Chân vuông góc lên đường thẳng & $H = A + \operatorname{proj}_{\mathbf{u}}(\overrightarrow{AM})$ \\
\hline
Khoảng cách điểm--đường thẳng & $d(M,d) = \bigl\|\mathbf{v} - \operatorname{proj}_{\mathbf{u}}(\mathbf{v})\bigr\|$ \\
\hline
Hình chiếu vectơ lên mặt phẳng & $\mathbf{w}_{\parallel} = \mathbf{w} - \operatorname{proj}_{\mathbf{n}}(\mathbf{w})$ \\
\hline
Khoảng cách điểm--mặt phẳng & $d(M,(P)) = \dfrac{|ax_0+by_0+cz_0+d|}{\sqrt{a^2+b^2+c^2}}$ \\
\hline
Hình chiếu điểm lên mặt phẳng & $H = M - \operatorname{proj}_{\mathbf{n}}(\overrightarrow{AM})$ \\
\hline
\end{tabular}
\end{center}

Sáu dòng chỉ là dòng đầu tiên, cộng thao tác chọn đúng cặp vectơ và ghim điểm gốc.

\section{Bài tập tự luyện (có đáp án)}

\textbf{Bài 1.} Cho $d:\; \dfrac{x+1}{1} = \dfrac{y-2}{2} = \dfrac{z}{-2}$ và $M(0;5;-1)$. Tìm hình chiếu vuông góc của $M$ lên $d$ và khoảng cách $d(M,d)$.

\textbf{Bài 2.} Cho $(P): x - 2y + 2z + 3 = 0$ và $M(1;1;1)$. Tìm hình chiếu của $M$ lên $(P)$ và khoảng cách từ $M$ đến $(P)$ bằng phương pháp hình chiếu.

\textbf{Bài 3 (nâng cao).} Cho đường thẳng $\Delta$ và hai điểm $A$, $B$ cùng phía so với $\Delta$. Tìm $H$ trên $\Delta$ sao cho $AH+HB$ nhỏ nhất. Gợi ý của tôi: lấy hình chiếu $B'$ của $B$ lên $\Delta$, đối xứng $B$ qua $B'$ được $B''$; giao $AB''$ với $\Delta$ chính là $H$.

\textbf{Đáp số.}

\begin{enumerate}
  \item $H = (0;4;-2)$; $d(M,d) = \sqrt{2}$.
  \item $H = \left(\dfrac{5}{9};\,\dfrac{17}{9};\,\dfrac{1}{9}\right)$; $d\bigl(M,(P)\bigr) = \dfrac{4}{3}$.
\end{enumerate}

\section{Điều tôi rút ra}

Trong đại số tuyến tính, $\operatorname{proj}_{\mathbf{u}}$ không chỉ là mẹo: nó là một \textbf{ánh xạ tuyến tính} --- biến mọi vectơ thành thành phần thẳng hàng với $\mathbf{u}$, bảo toàn tổ hợp tuyến tính, và là viên gạch nền của trực giao hóa Gram--Schmidt, hồi quy tuyến tính, phân tích dữ liệu sau này.

Còn với phòng thi: sáu dạng bài Oxyz tôi từng học thuộc riêng lẻ giờ nằm gọn trong một phép tách $\mathbf{v} = \operatorname{proj}_{\mathbf{u}}(\mathbf{v}) + \mathbf{v}_{\perp}$ --- vế trái trả lời ``cái gì nằm trên'', độ dài vế phải trả lời ``xa bao nhiêu''. Học thuộc sáu công thức thì quên một là mất một dạng; hiểu một ánh xạ tuyến tính thì quên cũng viết lại được trong ba dòng.

Đó là toàn bộ ý nghĩa của cái tên ``tà đạo'' mà tôi đặt cho chuỗi này.

\end{document}

